\begin{theorem}[The fixed 196-variable quartic]\label{q196:fixed}
Let $q_1^{[196]}(A),\ldots,q_{182}^{[196]}(A)$ be the integer
quadratic residuals printed in Section~\ref{q196:rows}. Set
\[
 N_{\rm BB}=\BB(2^{1000}),\qquad
 A_{\rm BB}=149^{2N_{\rm BB}+1},\qquad
 P_{\rm BB}^{[196]}=\sum_{i=1}^{182}\bigl(q_i^{[196]}(A_{\rm BB})\bigr)^2.
\]
Then, for every $x\in\Q$,
\[
 x\in\Z\quad\Longleftrightarrow\quad
 \exists(y_1,\ldots,y_{195})\in\Q^{195}:
 P_{\rm BB}^{[196]}(x,y_1,\ldots,y_{195})=0.
\]
The polynomial has degree four, content one and exactly 2943 nonzero
monomials. Its coefficients are integer polynomials of degree at most two
in the fixed integer $A_{\rm BB}$, and its largest absolute coefficient is
\[
 H(P_{\rm BB}^{[196]})=149^{4\BB(2^{1000})+2}<\BB(2^{2000}).
\]
The parameter $A_{\rm BB}$ is not an additional variable. The variable
count is an upper bound, not a minimum-variable claim.
\end{theorem}
